% ID: MAT-STA-GRA-001
% Titolo: Differenze e anomalie in una serie temporale
% Disciplina: Matematica
% Area: statistica
% Argomento: grafici_statistici
% Classe: Prima secondaria di secondo grado
% Difficolta: 3
% Tipo: laboratorio
% Tag: matematica, statistica, grafici_statistici
% Fonte: laboratorio_dati_approfondimento_anonimizzato
% Provenienza: materiale fornito dal docente; testo anonimizzato; soluzione redatta per Exergo

\begin{esercizio}
Ai tempi $0,2,4,6,8,10$ minuti sono associati i dati simulati $68,61,55,50,46,43$ gradi Celsius.
Calcola differenze successive e diminuzione media per minuto; traccia quest'ultima rispetto al tempo finale di ogni intervallo.
Classifica come ALTA una temperatura di almeno 50 e conta tali valori.
Ripeti l'ultimo intervallo sostituendo 43 con 51: che cosa segnala?
\end{esercizio}

\begin{soluzione}
Le differenze sono $-7,-6,-5,-4,-3$; dividendo gli opposti per 2 si ottengono
$3{,}5;3;2{,}5;2;1{,}5$ gradi al minuto, associati ai tempi $2,4,6,8,10$.
I valori almeno 50 sono quattro, incluso 50.
Con 51 l'ultima differenza è $51-46=+5$ e la diminuzione per minuto è $-2{,}5$:
l'ultimo intervallo descrive un aumento, quindi il dato va verificato.
Non è necessariamente un errore di formula. In Excel si possono usare
\texttt{=B6-B5}, \texttt{=-C6/(A6-A5)}, \texttt{=SE(B5>=50;"ALTA";"BASSA")}
e \texttt{=CONTA.SE(B5:B10;">=50")}.
\end{soluzione}
